\begin{center}
\textbf{RESUMO}
\end{center}

$\!$\\

\hspace{-1.3cm}\textbf{BANDEIRA}, Vinicius; \textbf{COUTO}, Pedro Henrique; \textbf{LIMA}, Caio Phelipe; \textbf{OLIVEIRA}, Iverson; \textbf{COSTA}, Rafael; \textbf{MUNIZ}, Rafael. \textit{Modelagem Dinâmica de um Robô Articulado de Baixo Custo para Locomoção em Corda}. 55 f. Relatório técnico (Graduação em Engenharia Mecânica) - Instituto Politécnico do Rio de Janeiro, Universidade do Estado do Rio de Janeiro~(UERJ), Nova Friburgo, 2026.

\vspace{.2cm}

Este relatório apresenta a modelagem dinâmica de um robô articulado de baixo custo, construído com palitos de picolé e de churrasco e acionado por um motor de corrente contínua com caixa de redução, capaz de se deslocar ao longo de uma corda. A massa, as dimensões e os ângulos de cada peça foram levantados com balança, paquímetro e fotografias, e o tensor de inércia de cada componente foi calculado a partir de sólidos homogêneos elementares e composto pelo teorema dos eixos paralelos, resultando em um robô de 103,1~g com centro de massa 21,4~mm abaixo do eixo do motor. O mecanismo, em que duas manivelas defasadas de 180$^\circ$ acionam braços que funcionam como alavancas em torno do eixo do quadril e terminam em ganchos, foi equacionado com as relações vetoriais de velocidade e aceleração do corpo rígido e com as equações de Newton-Euler projetadas nas velocidades generalizadas (método de Kane), verificadas contra a formulação de Lagrange com diferença relativa da ordem de $10^{-16}$. O contato dos ganchos com a corda foi tratado como um pino, com troca de mão por impacto plástico. A simulação com o motor a 30~rpm no sentido anti-horário mostrou o avanço do robô para a frente, com deslocamento do centro de massa de 131~mm em 8~s e velocidade média de cerca de 23~mm/s, contra 43,5~mm/s previstos pela cinemática pura, e com força vertical da corda sempre positiva. O avanço depende do amortecimento do balanço do chassi, que foi suposto, e da regra de engate adotada.

\vspace{1cm}

\hspace{-1.3cm}Palavras-chave: Dinâmica de corpos rígidos, Newton-Euler, Tensor de inércia, Robô articulado, Simulação.
